% f04_fabric variant c -- the fabric as a cycle chart: GENERATE, COMPUTE and
% SEND on one PE, the packet crossing the NoC, and the receiving PE releasing a
% COMPUTE, with every gap read off the RTL rather than assumed.
% Target: IEEE double column (figure*).  Fixed 16.9cm canvas -- the standalone
% wrapper is varwidth=17cm, so 17.6cm would be cropped in the guide preview.
% Data : rtl/noc_ni.sv (GCI state machine: gci_run/gci_ctr, GCI_LATENCY=4 ->
%          \GciLatency; gci_done is registered one cycle after gci_ctr reaches
%          GCI_LATENCY-1, and the DMEM write plus scoreboard set land one cycle
%          after that, so gci_start -> visible presence bit is \GciLatency+2 = 6);
%        rtl/btree_fsm_fast.sv (IF/DE/EX; if_advance = de_advance = ex_advance &&
%          !de_hazard; scoreboard_hazard, gci_data_hazard, ex_send_blocked;
%          piu_active accept-then-write, two cycles; combinational fp64_add;
%          TIMEOUT_MAX 10,000 and WATCHDOG_MAX -> \WatchdogMax);
%        rtl/link_tx.sv (credit-gated tx_valid), rtl/noc_router.sv;
%        data/claim_macros.tex, data/campaign_macros.tex.
% Strategy: TIME, not structure.  a traces one packet through the boxes and b
% partitions the invariants by layer; both are static, so neither can show the
% thing the prose spends most of its words on -- that a PE never waits on a
% clock, it waits on a presence bit, and that every latency in the machine is
% either a fixed RTL constant or explicitly topology-dependent.  Reading down a
% column gives the machine state in that cycle; reading across a row gives one
% unit's duty cycle.  d dissects the router instead.
% Colour: the instruction stream is the compliant datapath (spBlue -- pale while
% a stage merely holds, solid while it does work); a commit, meaning a DMEM
% write and the presence bit it sets, is confirmed (spGreen, heavy rule), and
% the bands where a presence bit stays 1 are the same green, tinted; every stall
% is the presence gate holding rather than a fault, so it is contended, not
% violating (spAmber).  The one duration the RTL does not fix is drawn on the
% blue path but hatched in spMute.  The two rows whose agents write DMEM with no
% instruction -- PIU and GCI -- carry spPurple labels, as in variants a and b.
% The fill ramp (held < issue < execute < stall < commit) and the rule weights
% carry the same ordering with the colour stripped.
\begin{tikzpicture}[
  font=\scriptsize,
  rowlbl/.style={font=\tiny, anchor=east, align=right, inner sep=0pt, text=spInk},
  grp/.style={font=\tiny\bfseries, anchor=west, inner sep=0pt, text=spInk},
  cyc/.style={font=\tiny, anchor=south, inner sep=0.5pt, text=spMute},
  note/.style={font=\tiny, inner sep=1pt, align=left, text=spInk},
  % the five block roles -- fill value AND rule weight order them without colour
  fHold/.style ={draw=spBlue!60, fill=spBlue!5,  line width=0.3pt,
                 dash pattern=on 1.3pt off 1.1pt},
  fIss/.style  ={draw=spBlue,    fill=spBlueF,   line width=0.3pt},
  fEx/.style   ={draw=spBlue,    fill=spBlue!22, line width=0.45pt},
  fStl/.style  ={draw=spAmber!85!spInk, fill=spAmber!55, line width=0.45pt},
  fCom/.style  ={draw=spGreen!80!spInk, fill=spGreen!60, line width=0.7pt},
  fPres/.style ={draw=spGreen,   fill=spGreenF,  line width=0.3pt},
]
\def\W{0.655}   % one cycle
\def\H{0.46}    % one row
\def\X{3.15}    % left edge of cycle 0
\def\Y{8.05}    % top edge of row 0

\newcommand{\cxx}[1]{\X+#1*\W}
\newcommand{\ryy}[1]{\Y-#1*\H}

% a filled block: {row}{first cycle}{span}{role style}{label}
\newcommand{\blk}[5]{%
  \filldraw[#4]
    (\cxx{#2}+0.04,\ryy{#1}-\H+0.05) rectangle (\cxx{#2}+#3*\W-0.04,\ryy{#1}-0.05);
  \node[font=\tiny, inner sep=0pt, text=spInk]
    at (\cxx{#2}+#3*\W/2,\ryy{#1}-\H/2) {#5};}
\newcommand{\hblk}[4]{%
  \filldraw[draw=spBlue, line width=0.3pt, pattern=north east lines,
            pattern color=spMute]
    (\cxx{#2}+0.04,\ryy{#1}-\H+0.05) rectangle (\cxx{#2}+#3*\W-0.04,\ryy{#1}-0.05);
  \node[font=\tiny, inner sep=1pt, fill=white, fill opacity=0.85, text opacity=1,
        text=spInk]
    at (\cxx{#2}+#3*\W/2,\ryy{#1}-\H/2) {#4};}

\path (0,0.42) rectangle (16.9,8.72);

% ---- cycle ruler -----------------------------------------------------------
\foreach \c in {0,...,19} {
  \draw[spGrid, very thin] (\cxx{\c},\ryy{13}) -- (\cxx{\c},\Y);
  \node[cyc] at (\cxx{\c}+\W/2,\Y+0.04) {\c};
}
\draw[spGrid, very thin] (\cxx{20},\ryy{13}) -- (\cxx{20},\Y);
\node[font=\tiny, anchor=east, text=spMute] at (\X-0.10,\Y+0.14) {cycle};
\draw[spMute, semithick] (\X,\Y) -- (\cxx{20},\Y);
\foreach \r in {5,8} \draw[spMute!60, thin] (0.05,\ryy{\r}) -- (\cxx{20},\ryy{\r});
\draw[spMute, semithick] (0.05,\ryy{13}) -- (\cxx{20},\ryy{13});

% =====================================================================
% BAND 1 -- the emitting PE, p
% =====================================================================
\node[grp, rotate=90, anchor=center, align=center] at (0.20,\ryy{3}) {PE $p$ \ \sv{btree\_fsm\_fast}};

\node[rowlbl] at (\X-0.10,\ryy{1}-\H/2) {\textsf{IF} IMEM read};
\blk{1}{0}{1}{fIss}{\texttt{i2}}
\blk{1}{1}{6}{fHold}{held}
\blk{1}{7}{1}{fIss}{---}

\node[rowlbl] at (\X-0.10,\ryy{2}-\H/2) {\textsf{DE} presence gate};
\blk{2}{0}{1}{fIss}{\texttt{i1}}
\blk{2}{1}{6}{fStl}{\sv{de\_hazard}}
\blk{2}{7}{1}{fIss}{\texttt{i2}}
\blk{2}{8}{1}{fStl}{}

\node[rowlbl] at (\X-0.10,\ryy{3}-\H/2) {\textsf{EX} add / TX enq.};
\blk{3}{0}{1}{fEx}{\texttt{i0}}
\blk{3}{7}{1}{fEx}{\texttt{i1}}
\blk{3}{9}{1}{fEx}{\texttt{i2}}

\node[rowlbl, text=spPurple!85!spInk] at (\X-0.10,\ryy{4}-\H/2) {GCI agent (\sv{noc\_ni})};
\blk{4}{1}{4}{fIss}{\sv{gci\_ctr}$=0..3$}
\blk{4}{5}{1}{fCom}{\sv{done}}

\node[rowlbl] at (\X-0.10,\ryy{5}-\H/2) {DMEM wr / \sv{scoreboard}};
\blk{5}{6}{14}{fPres}{\sv{scoreboard[$d$]}$=1$, monotone}
\blk{5}{5}{1}{fCom}{wr $d$}
\blk{5}{7}{1}{fCom}{wr $e$}

% the interval the figure is named for -- drawn inside the empty left half of row 5
\draw[{Stealth[length=1.6mm]}-{Stealth[length=1.6mm]}, semithick, spInk]
  (\cxx{0}+0.06,\ryy{5}-\H/2) -- (\cxx{5}-0.02,\ryy{5}-\H/2);
\node[note, fill=white, inner sep=0.8pt, anchor=center]
  at (\cxx{2.5},\ryy{5}-\H/2) {GCI gap $=\GciLatency{+}2=6$\,cy};

% =====================================================================
% BAND 2 -- the network
% =====================================================================
\node[grp, rotate=90, anchor=center] at (0.20,\ryy{6}-\H) {network};
\node[rowlbl] at (\X-0.10,\ryy{6}-\H/2) {TX $\to$ \sv{link\_tx}};
\blk{6}{9}{1}{fEx}{enq}
\blk{6}{10}{1}{fCom}{\sv{tx\_v}}
\node[note, anchor=west, text=spMute] at (\cxx{11}+0.10,\ryy{6}-\H/2)
  {sends only while \sv{credits}$>0$};

\node[rowlbl] at (\X-0.10,\ryy{7}-\H/2) {NI $+$ \sv{noc\_router}};
\hblk{7}{11}{4}{$\ge 1$ hop}

% =====================================================================
% BAND 3 -- the receiving PE, q
% =====================================================================
\node[grp, rotate=90, anchor=center] at (0.20,\ryy{10}-\H) {PE $q$ (receiver)};
\node[rowlbl, text=spPurple!85!spInk] at (\X-0.10,\ryy{9}-\H/2) {PIU (RDMA-push)};
\blk{9}{15}{1}{fIss}{acc}
\blk{9}{16}{1}{fCom}{wr}

\node[rowlbl] at (\X-0.10,\ryy{10}-\H/2) {\sv{scoreboard[$e'$]}};
\blk{10}{17}{3}{fPres}{$=1$}

\node[rowlbl] at (\X-0.10,\ryy{11}-\H/2) {\textsf{DE} at $q$};
\blk{11}{0}{17}{fStl}{a \texttt{COMPUTE} at $q$ stalled here the whole time --- waiting on a \emph{bit}, not a cycle count}
\blk{11}{17}{1}{fIss}{\texttt{go}}

\node[rowlbl] at (\X-0.10,\ryy{12}-\H/2) {\textsf{EX} at $q$};
\blk{12}{18}{1}{fEx}{\texttt{add}}
\draw[-{Stealth[length=1.6mm]}, semithick, spGreen]
  (\cxx{17.5},\ryy{10}-\H) -- (\cxx{17.5},\ryy{11});

% =====================================================================
% legend and the argument
% =====================================================================
\filldraw[fCom]
  (0.10,\ryy{13}-0.42) rectangle (0.42,\ryy{13}-0.16);
\node[note, anchor=west] at (0.48,\ryy{13}-0.29) {commit};
\filldraw[fStl]
  (1.55,\ryy{13}-0.42) rectangle (1.87,\ryy{13}-0.16);
\node[note, anchor=west] at (1.93,\ryy{13}-0.29) {stall (presence gate)};
\filldraw[fIss]
  (4.55,\ryy{13}-0.42) rectangle (4.87,\ryy{13}-0.16);
\node[note, anchor=west] at (4.93,\ryy{13}-0.29)
  {issue \quad $\texttt{i0}=$\texttt{GENERATE}\,$d$; \ $\texttt{i1}=$\texttt{COMPUTE}\,$e\!\gets\! d+\mathrm{DMEM}[0]$; \ $\texttt{i2}=$\texttt{SEND}\,$e\to q$ \quad (\texttt{DMEM[0]} is the identity, always present)};

\node[note, anchor=north west, text width=16.7cm] at (0.05,\ryy{13}-0.62)
  {\textbf{Reading the chart.} Every stall in it is a presence gate, never a
   timer: the compiler emits no delays and no barriers. The GCI gap is
   \GciLatency{} cycles counted in \sv{gci\_ctr}, $+1$ to register \sv{gci\_done},
   $+1$ for the DMEM write and the scoreboard set --- the only fixed non-unit
   latency inside a PE --- and cycles 1--6 are held twice over, by
   \sv{scoreboard[$d$]}$=0$ and again by \sv{gci\_data\_hazard}, which catches
   the in-flight case the scoreboard cannot see. The hatched row is the only
   quantity the figure declines to name, because it is set by the adjacency the
   compiler loaded rather than by the RTL. Changing either row changes
   \emph{when} the bottom row fires, never \emph{whether} it fires or what it
   computes --- which is why the root word is identical across all
   \CampPassing{} passing runs while their cycle counts are not. Should either
   stall fail to clear, \sv{stall\_timer} $\ge 10{,}000$ or the retirement
   watchdog $\ge \WatchdogMax$ ends the run in \sv{P\_ERROR} with a recorded
   \sv{error\_pc}, so the chart has no path to a wrong answer, only to a halt.};
\end{tikzpicture}
